% =====================================================================
% Présentation Beamer : « Pairs Trading with Wavelet Transform »
% Université Paris-Dauphine -- Master 272 -- Arbitrage Statistique
% Auteurs : Yassine Mannai, Antonin Bezard, Issam Fradi
% Compilation : pdflatex -interaction=nonstopmode presentation.tex
%          ou : .\compile.ps1
% =====================================================================
\documentclass[10pt,aspectratio=169]{beamer}

\usepackage[utf8]{inputenc}
\usepackage[T1]{fontenc}
\usepackage[french]{babel}
\usepackage{lmodern}
\usepackage{amsmath, amssymb}
\usepackage{booktabs}
\usepackage{array}
\usepackage{graphicx}
\usepackage{hyperref}

\usetheme{Madrid}
\usecolortheme{seahorse}
\setbeamertemplate{navigation symbols}{}
\setbeamertemplate{caption}[numbered]
\setbeamerfont{caption}{size=\footnotesize}
\setbeamersize{text margin left=12pt, text margin right=12pt}

\graphicspath{{../outputs/figures/}}

% Repli gracieux quand une figure n'est pas (encore) générée sur disque.
\newcommand{\figureorplaceholder}[2][height=0.72\textheight]{%
  \IfFileExists{../outputs/figures/#2.png}
    {\includegraphics[#1]{#2}}
    {\fbox{\parbox{0.42\textwidth}{\centering\footnotesize\ttfamily\detokenize{#2}.png\\[2pt]\upshape\itshape (lancer \texttt{py research/main.py -{}-full} pour g\'en\'erer)}}}%
}

\title[Pairs Trading par Ondelettes]{Pairs Trading par Transformée en Ondelettes\\\small{Réplication de Eroglu, Yener \& Yigit (2023)}}
\subtitle{Quantitative Finance, 23(7--8), 1129--1154}
\author{Yassine Mannai \and Antonin Bezard \and Issam Fradi}
\institute{Université Paris-Dauphine --- Master 272\\Arbitrage Statistique}
\date{\today}

\begin{document}

% ---------------------------------------------------------------------
\begin{frame}
  \titlepage
\end{frame}

% ---------------------------------------------------------------------
\begin{frame}{Plan}
  \tableofcontents
\end{frame}

% =====================================================================
\section{Le papier}
% =====================================================================
\begin{frame}{Le papier en une diapositive}
  \textbf{Eroglu, Yener \& Yigit (2023) --- « Pairs trading with wavelet transform », \textit{Quantitative Finance}.}
  \vskip4pt
  \begin{itemize}
    \item Le pairs trading classique repose sur un spread de prix qui mélange
          \emph{tendance} (marche aléatoire) et \emph{bruit haute fréquence}.
    \item Ce bruit gonfle la variance du spread, brouille les signaux
          d'entrée/sortie et déclenche des sorties prématurées.
    \item Idée : remplacer le prix brut par sa \textbf{composante lisse
          d'ondelettes} (MODWT niveau 1+, symlet 20) avant de former le
          spread. La partie lisse conserve la dynamique lente et
          écarte la contamination haute fréquence.
    \item Résultats : 11{,}82~\% (distance) et 9{,}66~\% (cointégration) de
          rendement moyen par période contre $-0{,}55~\%$ et $-1{,}81~\%$
          pour la variante standard. Sharpe : 3{,}69 / 2{,}82 contre
          $-0{,}21$ / $-0{,}40$.
    \item Univers : S\&P~500, point-in-time, 1986--2020, 415 tickers finaux.
  \end{itemize}
\end{frame}

% =====================================================================
\section{Stratégie}
% =====================================================================
\begin{frame}{Découpage formation / trading}
  \begin{itemize}
    \item L'historique est découpé en blocs consécutifs de 252 jours ouvrés.
    \item Pour chaque période $n$ :
      \begin{itemize}
        \item formation = bloc $k$ (in-sample),
        \item trading  = bloc $k+1$ (out-of-sample).
      \end{itemize}
    \item La sélection des paires se fait \textbf{uniquement} sur le bloc de
          formation, à partir de l'univers survivant à cette date
          (aucun look-ahead bias).
  \end{itemize}

  \vskip6pt
  \begin{center}
    \scriptsize
    \begin{tabular}{>{\bfseries}r p{0.65\linewidth}}
      \toprule
      Distance      & 20 paires minimisant $\sum (P_i^* - P_j^*)^2$ sur les prix normalisés. \\
      Cointégration & Test ADF d'Engle--Granger sur $\tilde{P}_i - \alpha - \beta \tilde{P}_j$ après pré-filtre par distance. \\
      \bottomrule
    \end{tabular}
  \end{center}
\end{frame}

\begin{frame}{L'étape ondelettes}
  \begin{itemize}
    \item Chaque série de prix est rééchelonnée, $\tilde{P}_t = P_t / P_0$.
    \item MODWT niveau 1+ avec le \textbf{symlet sym20}, condition aux
          bords périodique sur le bloc de trading (fidèle au papier,
          section 4.4).
    \item On conserve les \textbf{coefficients lisses} (passe-bas), on
          jette les détails (passe-haut), puis on reconstitue le prix
          synthétique.
    \item Le spread de la paire filtrée par ondelettes est alors
      \[ s_t = \tilde{P}_i^{\text{wav}} - \alpha - \beta\, \tilde{P}_j^{\text{wav}}, \]
      avec $(\alpha,\beta)$ estimés par MCO sur le bloc de formation.
    \item La variante standard utilise la même formule \emph{sans} l'étape
          MODWT --- c'est la seule différence entre les courbes de chaque figure.
  \end{itemize}
\end{frame}

\begin{frame}{Règle de trading}
  \begin{itemize}
    \item Calculer $\sigma$ du spread sur la fenêtre de formation.
    \item Entrer dès que $|s_t| > 2\sigma$ (vente du jambage haut, achat du jambage bas).
    \item Sortir quand $s_t$ recroise zéro pour la première fois.
    \item Clôture forcée en fin de bloc de trading ; la variante
          \textbf{forced-close-only} de la Table~22 n'utilise aucune autre sortie.
    \item Coûts de transaction : coût par action appliqué à l'entrée et
          à la sortie. Nous balayons $tc \in \{0~;~0{,}001~;~0{,}005\}$
          pour les Tables~18--21.
  \end{itemize}
\end{frame}

% =====================================================================
\section{Pipeline de réplication}
% =====================================================================
\begin{frame}{Architecture de la réplication}
  \footnotesize
  \begin{tabular}{>{\ttfamily}l p{0.65\linewidth}}
    \toprule
    research/main.py            & CLI unifiée. \texttt{--full} = run complet du papier, \texttt{--tc-sweep} = Tables~18--21. \\
    paper\_replication/core/    & Moteur par période : periods, selection, spread, trading, wavelet, metrics. \\
    paper\_replication/analytics/ & Régressions d'asset pricing (Fama-French, q-factor, ICAPM de Petkova). \\
    paper\_replication/outputs/ & 22 tables numérotées + 12 figures + résumé de validation. \\
    paper\_replication/bootstrap/ & Récupération Bloomberg et rafraîchissement des facteurs CSV. \\
    \bottomrule
  \end{tabular}
  \vskip6pt
  \begin{itemize}
    \item Polars + numpy + statsmodels (HAC Newey--West), export PNG via matplotlib.
    \item Sélection des paires et trading : Python pur, déterministe, sans MATLAB.
    \item 31 tests unitaires (\texttt{py -m pytest -q}) couvrant pipeline,
          tables et seuils de validation.
  \end{itemize}
\end{frame}

% =====================================================================
\section{Résultats principaux}
% =====================================================================
\begin{frame}{Chiffres-clés (Tables 4 \& 5)}
  \centering
  \footnotesize
  \begin{tabular}{l rr rr}
    \toprule
                  & \multicolumn{2}{c}{\bfseries Rendement moyen (\%)} & \multicolumn{2}{c}{\bfseries Sharpe} \\
    \cmidrule(lr){2-3}\cmidrule(lr){4-5}
    Méthode       & Standard & Ondelettes & Standard & Ondelettes \\
    \midrule
    Distance      & $-0{,}55$ & $11{,}82$ & $-0{,}21$ & $3{,}69$ \\
    Cointégration & $-1{,}81$ & $9{,}66$  & $-0{,}40$ & $2{,}82$ \\
    \bottomrule
  \end{tabular}
  \vskip6pt
  Les tolérances de réplication sont dans \texttt{paper\_validation\_summary.csv} :
  \texttt{PASS} si $|\Delta \text{mean}| \leq 1{,}0$~pp et
  $|\Delta \text{Sharpe}| \leq 0{,}30$.
\end{frame}

% =====================================================================
\section{Figures reproduites}
% =====================================================================
\newcommand{\figurepair}[2]{%
  \begin{minipage}[c]{0.48\textwidth}\centering\figureorplaceholder[width=\linewidth]{#1}\end{minipage}\hfill%
  \begin{minipage}[c]{0.48\textwidth}\centering\figureorplaceholder[width=\linewidth]{#2}\end{minipage}%
}

\begin{frame}{Figure 4 --- Rendements cumulés en excès}
  \figurepair{figure4_cumulative_excess_returns_distance}{figure4_cumulative_excess_returns_cointegration}
  \begin{center}\tiny Gauche : méthode des distances. Droite : méthode de cointégration.\end{center}
\end{frame}

\begin{frame}{Figure 5 --- Ratios de Sharpe quotidiens}
  \figurepair{figure5_daily_sharpe_ratios_distance}{figure5_daily_sharpe_ratios_cointegration}
\end{frame}

\begin{frame}{Figure 8 --- Corrélation du bruit filtré}
  \figurepair{figure8_filtered_noise_correlation_distance}{figure8_filtered_noise_correlation_cointegration}
\end{frame}

\begin{frame}{Figure 9 --- Rendements anormaux annuels}
  \figurepair{figure9_yearly_abnormal_returns_distance}{figure9_yearly_abnormal_returns_cointegration}
\end{frame}

% =====================================================================
\section{Bilan}
% =====================================================================
\begin{frame}{Couverture de la réplication}
  \footnotesize
  \begin{tabular}{l l p{0.55\linewidth}}
    \toprule
    Éléments              & Statut      & Notes \\
    \midrule
    Tables 1--9           & reproduites & run par défaut \\
    Table 10              & reproduite  & ACP, Python pur (auparavant marquée MATLAB-only) \\
    Tables 11, 13--15     & reproduites & nécessitent \texttt{research/data/factors.csv} \\
    Table 12              & reproduite  & Monte-Carlo contamination MODWT \\
    Tables 16--17         & reproduites & balayages classe d'ondelettes et horizon \\
    Tables 18--21         & opt-in      & \texttt{--tc-sweep 0.0 0.001 0.005} \\
    Table 22              & reproduite  & forced-close-only (paires standard) \\
    Figures 1--12         & reproduites & matplotlib, PNG style papier \\
    Résumé de validation  & livré       & \texttt{paper\_validation\_summary.csv} \\
    \bottomrule
  \end{tabular}
  \vskip8pt
  L'ensemble des 22 tables numérotées et des 12 figures est produit par une seule commande :
  \begin{center}
    \texttt{py research/main.py --full --tc-sweep 0.0 0.001 0.005}
  \end{center}
\end{frame}

\begin{frame}[plain]
  \begin{center}
    \Large\textbf{Merci de votre attention}\\[6pt]
    \small Code de réplication, slides et sorties :\\
    sous-arbre \texttt{research/} du backtesting framework.\\[8pt]
    \footnotesize Yassine Mannai \quad$\cdot$\quad Antonin Bezard \quad$\cdot$\quad Issam Fradi\\
    \tiny Université Paris-Dauphine --- Master 272 --- Arbitrage Statistique
  \end{center}
\end{frame}

\end{document}
